\documentclass[10pt,a4paper]{amsart}
\usepackage[margin=19mm]{geometry}
\usepackage{amsmath,amssymb,mathtools}
\usepackage[scr=rsfs,cal=euler]{mathalfa}
\usepackage{xfrac}
\usepackage[unicode,hidelinks,bookmarks=false]{hyperref}
\newcommand{\ie}{i.e.\ }
\newcommand{\cf}{cf.\ }
\newcommand{\resp}{respectively\ }
\newcommand{\Subsection}{\subsection}
\providecommand{\nobreakdash}{\nobreak\hbox{-}}
\setlength{\emergencystretch}{2em}
\multlinegap0pt
\usepackage{bm}
\usepackage{bbm}
\usepackage{dsfont}
\usepackage{xspace}
\usepackage{cancel}
\usepackage{mathtools}
\usepackage{amsthm}
\makeatletter
\@ifundefined{assumption}{\newtheorem{assumption}{Assumption}}{}
\@ifundefined{lemma}{\newtheorem{lemma}{Lemma}}{}
\@ifundefined{theorem}{\newtheorem{theorem}{Theorem}}{}
\makeatother
\newcommand{\bsu}{\boldsymbol{u}}
\newcommand{\bsx}{\boldsymbol{x}}
\newcommand{\itr}{\mathbb{I}}
\newcommand{\mix}{\mathrm{mix}}
\newcommand{\rd}{\,\mathrm{d}}
\newcommand{\real}{\mathbb{R}}
\newcommand{\tc}{\Tilde{c}}
\title[Quasi-Monte Carlo integration over $\mathbb{R}^s$ with bound]{Quasi-Monte Carlo integration over $\mathbb{R}^s$ with boundary-damping importance sampling}
\author{Zexin Pan, Du Ouyang, Zhijian He}
\date{Source: arXiv:2509.07509v2 (CC BY 4.0)}
\begin{document}
\maketitle

\noindent\emph{Source and licence.} Quasi-Monte Carlo integration over $\mathbb{R}^s$ with boundary-damping importance sampling. Version arXiv:2509.07509v2 (CC BY 4.0), distributed under the Creative Commons Attribution 4.0 International licence, as recorded in the arXiv licence metadata for this exact version and shown on the abstract page https://arxiv.org/abs/2509.07509v2. Full rights notice: licenses/arxiv-2509.07509-CC-BY-4.0.txt.\par\smallskip

\noindent\textit{Editorial scope.} Counted unit: the Theorem whose source label is \texttt{thm:Wqcase} in the article (source0.tex lines 600--604, source label \texttt{thm:Wqcase}), delivered together with the article's own complete original proof of it (source0.tex lines 605--642, one \texttt{proof} environment of 38 lines). No mathematics is written, repaired or re-derived: statement and proof are the article's own text line for line. Required context retained uncounted: assump:rho (lines 265--268); lem:dfw (lines 539--542); lem:rhoTbound (lines 557--562); lem:dwjbound (lines 579--584). Every definition, display and label of those lines is retained unchanged. Omitted from this package and not redistributed: 1--264, 269--538, 543--556, 563--578, 585--599, 643--1100. Nothing inside a retained excerpt is omitted or altered: every retained body is delivered exactly as the article's lines are written, so no omission receipt of any kind accompanies this package. Citations: the retained excerpts carry no citation command at all, so no bibliography entry of the article is transcribed and no reference is redistributed. Reference adaptation: none was needed and none was made; every reference inside the delivered unit resolves inside it, so no unresolved reference or citation is produced. Printed slips kept literal: no printed word, symbol, hypothesis or number of the counted statement or of its proof was altered. Layout and class: the article's own class is \texttt{siamart250211} with packages including lipsum, amsfonts, graphicx, tabularx, epstopdf, algorithmic, xcolor, subcaption, amsopn; the wrapper is the lane's pinned \texttt{amsart} editorial wrapper at 10pt on A4, which restates the theorem-like environments the retained text needs and adds the article's own macro definitions that the retained text needs (\texttt{\textbackslash bsu}, \texttt{\textbackslash bsx}, \texttt{\textbackslash itr}, \texttt{\textbackslash mix}, \texttt{\textbackslash rd}, \texttt{\textbackslash real}, \texttt{\textbackslash tc}), with extra wrapper packages (bm, bbm, dsfont, xspace, cancel, mathtools). The delivered numbering is the unit's own. Assets: no retained span contains a figure environment, a graphics-inclusion command, a PDF-page inclusion, a table or a TikZ picture; no image file is copied into this package, the package declares no asset dependency and no image file is redistributed with it. Licence: version arXiv:2509.07509v2 of this article is distributed under the Creative Commons Attribution 4.0 International licence, as recorded in the arXiv licence metadata for this exact version and shown on the abstract page https://arxiv.org/abs/2509.07509v2. A full rights notice is delivered at licenses/arxiv-2509.07509-CC-BY-4.0.txt. Characters: every retained excerpt is pure ASCII, so the delivered engine, XeLaTeX with its default Latin Modern text font, reports no missing character.
\begin{assumption}~\label{assump:rho}
Assume that $\varphi(x)>0$, $\varphi(x)=\varphi(-x)$, $\varphi_\infty:=\sup_{x\in \real} \varphi(x)<\infty$, and for any $\varepsilon>0$, there exists a constant $c_\varepsilon>0$ such that
$\varphi(x)\geq c_\varepsilon \Phi(x)^{1+\varepsilon}$ for any $x\le 0$. 
\end{assumption}

\begin{lemma}\label{lem:dfw}
    For $f\in W^{1,q}_{\mix}(\real^s,\varphi)$ and $f^w(\bsu)=w(\bsu)f\circ T(\bsu)$,
    $$\partial^{1{:}s} f^w (\bsu) =\sum_{v\subseteq 1{:}s}(\partial^{v} f) \circ T(\bsu) \prod_{j\in v} \frac{w^2_j(u_j)}{\varphi\circ T_j(u_j)}\prod_{j\in v^c} w'_j(u_j).$$
\end{lemma}

\begin{lemma}\label{lem:rhoTbound}
    Suppose $\varphi$ satisfies Assumption~\ref{assump:rho}. For $\varepsilon>0$ and $u\in (0,1/2]$,
    $$\varphi\circ T_j(u)\geq 
       c_{\varepsilon,p} \Big(  \theta_jw_{j}(u)\min\big((2u/\theta_j)^{p+1},1 \big)\Big)^{1+\varepsilon},$$
    where $c_{\varepsilon,p}$ are constants depending on $\varepsilon$ and $p$.
\end{lemma}

\begin{lemma}\label{lem:dwjbound}
    For $u\in (0,1/2]$,
    $$|w'_j(u)|\leq 
      c_p  \theta^{-1}_j w_j(u) \max\big( (2u/\theta_j)^{-p-1},1\big),$$
    where $c_{p}$ are constants depending on $p$.
\end{lemma}

\begin{theorem}\label{thm:Wqcase}
    For $f\in W^{1,q}_{\mix}(\real^s,\varphi)$ with $q>1$, $\varphi$ satisfying Assumption~\ref{assump:rho} and $\varepsilon\in (0,1-1/q)$, we have
    $$\Vert\partial^{1{:}s} f^w \Vert_{L^q(\itr^s)}\leq  \Big(\prod_{j=1}^sC_{\varepsilon,p,q}\theta^{-1-\varepsilon}_j\Big)\Vert f \Vert_{W^{1,q}_{\mix}(\real^s,\varphi)},$$
    where $C_{\varepsilon,p,q}$ is a constant depending on $\varepsilon,p$ and $q$.
\end{theorem}

\begin{proof}
    First we use Lemma~\ref{lem:dfw} to bound
    \begin{align}\label{eqn:partialfbound}
        |\partial^{1{:}s} f^w (\bsu)|\leq \sum_{v\subseteq 1{:}s}|(\partial^{v} f) \circ T(\bsu)| \prod_{j=1}^s w_j(u_j)^{1/q}\prod_{j\in v} \frac{w_j(u_j)^{2-1/q}}{\varphi\circ T_j(u_j)}\prod_{j\in v^c} \frac{|w'_j(u_j)|}{w_j(u_j)^{1/q}}.
    \end{align}
    By Lemma~\ref{lem:rhoTbound} and symmetry,
    \begin{align*}
        \sup_{u\in \itr }\frac{w_j(u)^{2-1/q}}{\varphi\circ T_j(u)}=&\sup_{u\in (0,1/2] }\frac{w_j(u)^{2-1/q}}{\varphi\circ T_j(u)}\\
        \leq &
       \tc^{-1}_{\varepsilon,p}  \theta^{-1-\varepsilon}_j\sup_{u\in (0,1/2] }w_{j}(u)^{1-1/q-\varepsilon}\max\big((2u/\theta_j)^{-(1+\varepsilon)(1+p)},1\big).
    \end{align*}
   Note that
   \begin{align*}
       &\sup_{u\in (0,1/2] }w_{j}(u)^{1-1/q-\varepsilon}\max\big((2u/\theta_j)^{-(1+\varepsilon)(1+p)},1\big)\\
       \leq & \max\Big(\sup_{u\in (0,\theta_j/2)}(2u/\theta_j)^{-(2-1/q)(1+p)}\exp\Big((1-1/q-\varepsilon)(2^p-(u/\theta_j)^{-p})\Big), 2\Big),
   \end{align*}
   which can be further bounded in terms of $\varepsilon,p$ and $q$. Hence
   \begin{equation}\label{eqn:woverphi}
      \sup_{u\in \itr }\frac{w_j(u)^{2-1/q}}{\varphi\circ T_j(u)}\leq C_1\theta^{-1-\varepsilon}_j 
   \end{equation}
   for $C_1$ depending on $\varepsilon,p$ and $q$. A similar calculation using Lemma~\ref{lem:dwjbound} shows
   \begin{equation}\label{eqn:dw}
       \sup_{u\in \itr } \frac{|w'_j(u_j)|}{w_j(u_j)^{1/q}} \leq C_2 \theta^{-1}_j
   \end{equation}
   for $C_2$ depending on $p$ and $q$. Using the above bounds, Equation~\eqref{eqn:partialfbound} becomes
   \begin{align*}
        |\partial^{1{:}s} f^w (\bsu)|\leq &\sum_{v\subseteq 1{:}s}|(\partial^{v} f) \circ T(\bsu)| \prod_{j=1}^s w_j(u_j)^{1/q} \prod_{j\in v} C_1\theta^{-1-\varepsilon}_j\prod_{j\in v^c} C_2\theta^{-1}_j\\
        \leq &  \Big(\prod_{j=1}^sC_{\varepsilon,p,q}\theta^{-1-\varepsilon}_j\Big) \Big(\sum_{v\subseteq 1{:}s}|(\partial^{v} f) \circ T(\bsu)|^q \prod_{j=1}^s w_j(u_j)\Big)^{1/q},
    \end{align*}
    where
    $$C_{\varepsilon,p,q} =\Big(\sum_{v\subseteq 1{:}s} C^{\frac{q}{q-1}|v|}_1C^{\frac{q}{q-1}(s-|v|)}_2\Big)^{\frac{1}{s}(1-\frac{1}{q})}=\Big(C^{\frac{q}{q-1}}_1+C^{\frac{q}{q-1}}_2\Big)^{1-\frac{1}{q}}.$$
    The conclusion then follows because
    \begin{align*}
     &\sum_{v\subseteq 1{:}s}\int_{\itr^s}|(\partial^{v} f) \circ T(\bsu)|^q \prod_{j=1}^s w_j(u_j)\rd \bsu\\
     =&\sum_{v\subseteq 1{:}s}\int_{\itr^s}|(\partial^{v} f) \circ T(\bsu)|^q \prod_{j=1}^s \varphi\circ T_j(u_j)T'_j(u_j)\rd \bsu\\
     = &\sum_{v\subseteq 1{:}s}\int_{\real^s}|\partial^{v} f(\bsx)|^q \prod_{j=1}^s \varphi(x_j)\rd \bsx=\Vert f \Vert^q_{W^{1,q}_{\mix}(\real^s,\varphi)}.
    \end{align*}
\end{proof}

\end{document}
